% Options for packages loaded elsewhere
\PassOptionsToPackage{unicode}{hyperref}
\PassOptionsToPackage{hyphens}{url}
\documentclass[
  american,
  11pt,
]{article}
\usepackage{xcolor}
\usepackage[margin=25mm]{geometry}
\usepackage{amsmath,amssymb}
\setcounter{secnumdepth}{-\maxdimen} % remove section numbering
\usepackage{iftex}
\ifPDFTeX
  \usepackage[T1]{fontenc}
  \usepackage[utf8]{inputenc}
  \usepackage{textcomp} % provide euro and other symbols
\else % if luatex or xetex
  \usepackage{unicode-math} % this also loads fontspec
  \defaultfontfeatures{Scale=MatchLowercase}
  \defaultfontfeatures[\rmfamily]{Ligatures=TeX,Scale=1}
\fi
\usepackage{lmodern}
\ifPDFTeX\else
  % xetex/luatex font selection
\fi
% Use upquote if available, for straight quotes in verbatim environments
\IfFileExists{upquote.sty}{\usepackage{upquote}}{}
\IfFileExists{microtype.sty}{% use microtype if available
  \usepackage[]{microtype}
  \UseMicrotypeSet[protrusion]{basicmath} % disable protrusion for tt fonts
}{}
\makeatletter
\@ifundefined{KOMAClassName}{% if non-KOMA class
  \IfFileExists{parskip.sty}{%
    \usepackage{parskip}
  }{% else
    \setlength{\parindent}{0pt}
    \setlength{\parskip}{6pt plus 2pt minus 1pt}}
}{% if KOMA class
  \KOMAoptions{parskip=half}}
\makeatother
% definitions for citeproc citations
\NewDocumentCommand\citeproctext{}{}
\NewDocumentCommand\citeproc{mm}{%
  \begingroup\def\citeproctext{#2}\cite{#1}\endgroup}
\makeatletter
 % allow citations to break across lines
 \let\@cite@ofmt\@firstofone
 % avoid brackets around text for \cite:
 \def\@biblabel#1{}
 \def\@cite#1#2{{#1\if@tempswa , #2\fi}}
\makeatother
\newlength{\cslhangindent}
\setlength{\cslhangindent}{1.5em}
\newlength{\csllabelwidth}
\setlength{\csllabelwidth}{3em}
\newenvironment{CSLReferences}[2] % #1 hanging-indent, #2 entry-spacing
 {\begin{list}{}{%
  \setlength{\itemindent}{0pt}
  \setlength{\leftmargin}{0pt}
  \setlength{\parsep}{0pt}
  % turn on hanging indent if param 1 is 1
  \ifodd #1
   \setlength{\leftmargin}{\cslhangindent}
   \setlength{\itemindent}{-1\cslhangindent}
  \fi
  % set entry spacing
  \setlength{\itemsep}{#2\baselineskip}}}
 {\end{list}}
\usepackage{calc}
\newcommand{\CSLBlock}[1]{\hfill\break\parbox[t]{\linewidth}{\strut\ignorespaces#1\strut}}
\newcommand{\CSLLeftMargin}[1]{\parbox[t]{\csllabelwidth}{\strut#1\strut}}
\newcommand{\CSLRightInline}[1]{\parbox[t]{\linewidth - \csllabelwidth}{\strut#1\strut}}
\newcommand{\CSLIndent}[1]{\hspace{\cslhangindent}#1}
\ifLuaTeX
\usepackage[bidi=basic,shorthands=off]{babel}
\else
\usepackage[bidi=default,shorthands=off]{babel}
\fi
\ifLuaTeX
  \usepackage{selnolig} % disable illegal ligatures
\fi
\setlength{\emergencystretch}{3em} % prevent overfull lines
\providecommand{\tightlist}{%
  \setlength{\itemsep}{0pt}\setlength{\parskip}{0pt}}
\usepackage{bookmark}
\IfFileExists{xurl.sty}{\usepackage{xurl}}{} % add URL line breaks if available
\urlstyle{same}
\hypersetup{
  pdftitle={Boundary-free weak influences: range closure{,} sharp information and feasible GLS},
  pdflang={en-US},
  hidelinks,
  pdfcreator={LaTeX via pandoc}}

\title{Boundary-free weak influences: range closure, sharp information
and feasible GLS}
\author{}
\date{}

\begin{document}
\maketitle

\subsection{1. Spaces and adjoint range}\label{spaces-and-adjoint-range}

Let \(J\in C^1([0,T];\mathbb R^{d\times d})\), and let the invertible
fundamental matrix satisfy \(Z'=JZ\), \(Z(0)=I_d\). Put \(A h=h'-Jh\)
and \(A^*\phi=-\phi'-J^T\phi\). Sensor coordinates have independent
stationary OU covariance with fixed positive scales. Its Cameron--Martin
space \(\mathcal H\) is \(H^1\) with the OU inner product. Write
\(\mathcal E=\mathcal H'\) for its continuous dual, equipped with the
covariance norm

\[
\|\lambda\|_C^2=\langle\lambda,C\lambda\rangle,
\qquad
\langle C^{-1}g,h\rangle=\langle g,h\rangle_{\mathcal H}.
\]

This definition includes ordinary integrated influences and endpoint
evaluation measures. It is the variance norm of the associated Gaussian
linear observation; the noise path itself need not belong to
\(\mathcal H\).

\textbf{Proposition 2 (exact range and closure).} In \(L^2\),

\[
A^*H_0^1
=\left\{\ell\in L^2:\int_0^T Z(t)^T\ell(t)dt=0\right\}.
\]

The range of compactly supported smooth tests is dense in this space in
\(L^2\). Its closure in \(\mathcal E\) is exactly

\[
\mathcal M=\overline{A^*C_c^\infty}^{\|\cdot\|_C}
=\{\lambda\in\mathcal E:\langle\lambda,Z_j\rangle=0, j=1,\ldots,d\}.
\]

Consequently \(\mathcal M\) has codimension \(d\).

\textbf{Proof.} Since \(Z^TA^*\phi=-(Z^T\phi)'\), zero endpoint values
imply the moment constraint. Conversely, for an \(L^2\) influence with
zero moment,

\[
\phi(t)=-Z(t)^{-T}\int_0^t Z(s)^T\ell(s)ds
\]

belongs to \(H_0^1\) and satisfies \(A^*\phi=\ell\). Approximation in
\(H_0^1\) by \(C_c^\infty\) and continuity of \(A^*:H_0^1\to L^2\) give
\(L^2\) density. For the stronger closure statement, a function
\(h\in\mathcal H\) annihilates every \(A^*\phi\) if and only if \(Ah=0\)
distributionally, hence \(h=Zc\). The Hahn--Banach annihilator
characterization of a closed subspace then gives the stated closure in
\(\mathcal H'\). The columns of \(Z\) are linearly independent elements
of \(\mathcal H\), so their dual constraints have rank \(d\).
\(\square\)

Zero values at the endpoints do not imply compact support. Equality for
literal compact tests, without a closure or an \(H_0^1\) qualification,
would be false. The relevant closure uses the covariance norm, which is
equivalent to the dual Sobolev norm for this positive OU model.

\subsection{2. Sharp covariance
reference}\label{sharp-covariance-reference}

Let \(S=[S_1,\ldots,S_p]\in\mathcal H^p\) be coefficient sensitivities,
and form the joint information blocks for \(S\) and \(Z\). Assume
\(\mathcal I_{uu}\) and the Schur complement are positive definite. Set

\[
R=S-Z\mathcal I_{uu}^{-1}\mathcal I_{u\theta},
\qquad
\mathcal I_e=\mathcal I_{\theta\theta}
-\mathcal I_{\theta u}\mathcal I_{uu}^{-1}\mathcal I_{u\theta}.
\]

Projection and best linear unbiased estimation in correlated continuous
regression are classical {[}1{]}. Here the adjoint-range
characterization identifies exactly which dynamical weak influences obey
that projection.

\textbf{Theorem 2 (sharp boundary-free influence covariance).} Let
\(\Lambda=[\lambda_1,\ldots,\lambda_p]\in\mathcal M^p\) satisfy
\(\langle\lambda_r,S_s\rangle=\delta_{rs}\). Its Gaussian influence
covariance obeys

\[
V(\Lambda)\succeq\mathcal I_e^{-1}.
\]

The bound is attained in \(\mathcal M\) by the matrix-valued dual
influence

\[
\Lambda_*=C^{-1}R\,\mathcal I_e^{-1}.
\]

It is the infimum for normalized compact-test influences. It is a
minimum among literal compact tests only when \(\Lambda_*\) belongs to
their adjoint range.

\textbf{Proof.} The \(\mathcal H\) projection makes
\(\langle C^{-1}R,Z\rangle=0\) and
\(\langle C^{-1}R,S\rangle=\mathcal I_e\). Thus \(\Lambda_*\) meets both
constraints and has covariance \(\mathcal I_e^{-1}\). The difference
\(\Delta=\Lambda-\Lambda_*\) pairs to zero with both \(S\) and \(Z\),
hence with \(R\). Covariance cross terms vanish and

\[
V(\Lambda)=\mathcal I_e^{-1}+V(\Delta).
\]

For attainability as an infimum, approximate each column of
\(\Lambda_*\) by compact-test influences in covariance norm. Their
normalization matrices converge to \(I_p\) by continuous pairing with
\(S\); right-multiplication by their inverse transpose enforces exact
normalization while preserving the range and covariance convergence.
\(\square\)

This is a matrix statement. Dividing a vector influence by a scalar
information is valid only when \(p=1\). The information loss is exactly
\(\mathcal I_{\theta u}\mathcal I_{uu}^{-1}\mathcal I_{u\theta}\);
covariance inflation is
\(\mathcal I_e^{-1}-\mathcal I_{\theta\theta}^{-1}\). There is no
direction-independent scalar loss percentage in general.

\subsection{3. Explicit OU
representative}\label{explicit-ou-representative}

For a scalar coordinate with \(g\in H^2\), integration by parts in the
stationary OU norm gives the dual measure

\[
C^{-1}g=
\frac{g-\tau^2g''}{2\sigma^2\tau}\,dt
+\frac{g(0)-\tau g'(0)}{2\sigma^2}\delta_0
+\frac{g(T)+\tau g'(T)}{2\sigma^2}\delta_T.
\]

For independent coordinates the formula applies coordinatewise, using
their respective scales. If only \(g\in H^1\), the Riesz representative
remains well-defined in \(\mathcal H'\) but need not be a density plus
finite endpoint measures.

For a smooth optimal representative write
\(\Lambda_*=m(t)dt+m_0\delta_0+m_T\delta_T\). On the open interval an
adjoint test matrix is

\[
\Phi_*(t)=-Z(t)^{-T}
\left[Z(0)^Tm_0+\int_0^t Z(s)^Tm(s)ds\right].
\]

Extend it by zero outside the observation interval. Its endpoint jumps
supply \(m_0=-\Phi_*(0+)\) and \(m_T=\Phi_*(T-)\), and distributionally
\(A^*\Phi_*=\Lambda_*\). These formulas include the initial atom;
omitting it in the cumulative integral is incorrect. The annihilation
condition supplies the terminal jump. Cutoff and smoothing approximate
this generalized test in covariance norm. These measures describe a
statistical reference; the implemented estimators keep compact tests and
add no explicit boundary observation terms.

For scalar exponential dynamics, \(z=e^{at}\), \(s=x_0te^{at}\),
\(k=\mathcal I_{au}/\mathcal I_{uu}\) and \(g=e^{at}(x_0t-k)\). Define

\[
A_0=\frac{(1-\tau^2a^2)x_0}{2\sigma^2\tau},\qquad
B_0=\frac{-(1-\tau^2a^2)k-2\tau^2ax_0}{2\sigma^2\tau},
\]

\[
q_0=\frac{-k-\tau(x_0-ak)}{2\sigma^2},\qquad
F_j(t)=\int_0^t s^je^{2as}ds.
\]

Then

\[
\phi_*(t)=-\frac{e^{-at}}{\mathcal I_e}
[q_0+A_0F_1(t)+B_0F_0(t)].
\]

For \(a\ne0\), \(F_0=\operatorname{expm1}(2at)/(2a)\) and
\(F_1=[e^{2at}(2at-1)+1]/(4a^2)\); their continuous limits at \(a=0\)
are \(t\) and \(t^2/2\). A pilot coefficient substitutes into this
expression before a compact cutoff is applied. The explicit multivariate
fundamental-matrix expression is an analytical identity, not a guarantee
of stable numerical inversion of \(Z(t)\); strongly contracting dynamics
can make that inversion ill-conditioned.

\subsection{4. GLS and the order of
limits}\label{gls-and-the-order-of-limits}

For tests \(\phi_k\), put \(\lambda_k=A^*\phi_k\),
\(D_{kr}=\langle\lambda_k,S_r\rangle\) and
\(\Sigma_{k\ell}=\langle\lambda_k,C\lambda_\ell\rangle\). Remove exact
linear dependencies. If \(D\) has full column rank, covariance-weighted
weak regression has local influence and covariance

\[
P=(D^T\Sigma^{-1}D)^{-1}D^T\Sigma^{-1},\qquad
V_K=(D^T\Sigma^{-1}D)^{-1}.
\]

WENDy's covariance iteration motivates this construction {[}2{]}, but
the released iid covariance is not the correct physical-OU covariance.
Under polynomial features Proposition 3 supplies the complete coupled
residual weights. Its order-one term determines the small-noise
efficient influence; higher orders correct residual uncertainty at
finite amplitude without by themselves removing coefficient-estimator
bias.

\textbf{Corollary 1 (small-noise, dense-span efficiency).} Let the noise
be \(\epsilon Z_C\) with known covariance shape \(C\). Assume a smooth
identifiable parameterization, a generating parameter interior to a
compact parameter set, consistent pilots, consistent covariance-shape
weights and quadrature error \(o(\epsilon)\). For each fixed finite test
family assume a regular local weak solution with the GLS linearization.
For the covariance conclusion additionally assume uniform integrability
of \(\|\epsilon^{-1}(\widehat\theta_{K,\epsilon}-\theta)\|^2\). A
sufficient implementation uses a uniformly stable smooth local solve,
Gaussian moment control on the local Taylor remainder, and bounded
fallback outside the regular region whose probability is
\(o(\epsilon^2)\). Then

\[
\epsilon^{-1}(\widehat\theta_{K,\epsilon}-\theta)
\Rightarrow N(0,V_K),\qquad
\epsilon^{-2}\operatorname{Cov}(\widehat\theta_{K,\epsilon})\to V_K.
\]

If nested test-influence spans are dense in \(\mathcal M\), then

\[
V_K\longrightarrow\mathcal I_{e,C}^{-1}.
\]

\textbf{Proof.} For fixed \(K\), a first-order expansion of the weak
equations gives \(D(\widehat\theta-\theta)=\epsilon r_1+o_p(\epsilon)\)
with the GLS map above. Consistent pilots and weights change the leading
map by \(o_p(1)\); derivatives of pilot-dependent tests multiply the
zero clean residual. Gaussianity gives the distributional assertion; the
stated risk condition justifies covariance convergence. In the Gaussian
dual space, finite-family information is the Gram matrix of the
orthogonal projection of \(C^{-1}R\) onto its test span. Density gives
projection convergence, and nested spans give increasing information;
continuity of inversion at positive definite \(\mathcal I_e\) gives
covariance convergence. \(\square\)

The conclusion is sequential: first small noise and controlled
quadrature for fixed \(K\), then dense spans. A diagonal joint limit
requires uniform remainder, conditioning and covariance-weight control
that are not supplied here. This theorem concerns regular WENDy-like
weak estimation with correct covariance and supplied identifiable
structure, not all implementations called WENDy, finite-noise sparse
selection, or a universal plateau for misspecified iid weighting. A
fixed basis cutoff must be compatible with the density assumption;
discarding an informative direction permanently can invalidate it.

\subsection*{References}\label{references}
\addcontentsline{toc}{subsection}{References}

\protect\phantomsection\label{refs}
\begin{CSLReferences}{0}{0}
\bibitem[\citeproctext]{ref-dette2016}
\CSLLeftMargin{{[}1{]} }%
\CSLRightInline{Dette H, Pepelyshev A, Zhigljavsky A. Optimal designs in
regression with correlated errors. \emph{The Annals of Statistics} 44
(2016) 113--152. \url{https://doi.org/10.1214/15-AOS1361}.}

\bibitem[\citeproctext]{ref-bortz2023}
\CSLLeftMargin{{[}2{]} }%
\CSLRightInline{Bortz DM, Messenger DA, Dukic V. Direct estimation of
parameters in {ODE} models using {WENDy}: Weak-form estimation of
nonlinear dynamics. \emph{Bulletin of Mathematical Biology} 85 (2023)
110. \url{https://doi.org/10.1007/s11538-023-01208-6}.}

\end{CSLReferences}

\end{document}
